\documentclass[12pt]{article}
\usepackage{amsmath,amssymb,amsfonts,amsthm}
\usepackage[margin=1in]{geometry}

\begin{document}

\begin{center}
{\Large\bfseries Chapter 9, Problem 5}\\[6pt]
Byron \& Fuller, \textit{Mathematics of Classical and Quantum Physics}
\end{center}

\bigskip

\textbf{Problem.} Consider the Hermitian integral operator $A$, whose square-integrable kernel is $a(s,t)$, and assume that all eigenvalues of $A$ are positive. Write a ``trial eigenvector'' $\varphi = \sum_{n=1}^{\nu} a_n \phi_n$ where $\{\phi_n\}$ is a given set of functions, and extremize $(\varphi, A\varphi)$ subject to $(\varphi, \varphi) = 1$.

(a) Show the resulting equation is $M\boldsymbol{\alpha} = \lambda N\boldsymbol{\alpha}$, where $M_{ij} = \int\!\!\int \phi_i^*(s)\,a(s,t)\,\phi_j(t)\,ds\,dt$ and $N_{ij} = \int \phi_i^*(s)\,\phi_j(s)\,ds$.

(b) Show $M$ and $N$ are Hermitian.

(c) Show the values of $\lambda$ are real.

(d) Discuss orthogonality of eigenvectors.

(e) Relate the largest $\lambda$ to the largest eigenvalue of $A$.

\bigskip
\textbf{Solution.}

\subsection*{(a) Derivation of the generalized eigenvalue equation}

We wish to extremize $(\varphi, A\varphi)$ subject to the constraint $(\varphi, \varphi) = 1$. Using the method of Lagrange multipliers, we extremize:
\[
\Phi[\{a_m\}] = (\varphi, A\varphi) - \lambda\bigl[(\varphi, \varphi) - 1\bigr].
\]

Computing each inner product:
\[
(\varphi, A\varphi) = \sum_{i,j} a_i^* a_j \int\!\!\int \phi_i^*(s)\,a(s,t)\,\phi_j(t)\,ds\,dt = \sum_{i,j} a_i^* M_{ij}\,a_j = \boldsymbol{\alpha}^\dagger M \boldsymbol{\alpha},
\]
\[
(\varphi, \varphi) = \sum_{i,j} a_i^* a_j \int \phi_i^*(s)\,\phi_j(s)\,ds = \sum_{i,j} a_i^* N_{ij}\,a_j = \boldsymbol{\alpha}^\dagger N \boldsymbol{\alpha}.
\]

Setting $\partial \Phi / \partial a_m^* = 0$ for each $m$:
\[
\sum_j M_{mj}\,a_j - \lambda \sum_j N_{mj}\,a_j = 0 \quad \Longrightarrow \quad \boxed{M\boldsymbol{\alpha} = \lambda N\boldsymbol{\alpha}}.
\]

\subsection*{(b) $M$ and $N$ are Hermitian}

For $M$:
\[
M_{ij}^* = \left(\int\!\!\int \phi_i^*(s)\,a(s,t)\,\phi_j(t)\,ds\,dt\right)^* = \int\!\!\int \phi_i(s)\,a^*(s,t)\,\phi_j^*(t)\,ds\,dt.
\]
Since $A$ is Hermitian, $a^*(s,t) = a(t,s)$, so:
\[
M_{ij}^* = \int\!\!\int \phi_j^*(t)\,a(t,s)\,\phi_i(s)\,dt\,ds = M_{ji}.
\]
Therefore $M^\dagger = M$, i.e., $M$ is Hermitian.

For $N$:
\[
N_{ij}^* = \left(\int \phi_i^*(s)\,\phi_j(s)\,ds\right)^* = \int \phi_i(s)\,\phi_j^*(s)\,ds = N_{ji}.
\]
So $N$ is also Hermitian. \hfill $\square$

\subsection*{(c) Eigenvalues $\lambda$ are real}

From $M\boldsymbol{\alpha} = \lambda N\boldsymbol{\alpha}$, take the inner product with $\boldsymbol{\alpha}$:
\[
\boldsymbol{\alpha}^\dagger M \boldsymbol{\alpha} = \lambda\,\boldsymbol{\alpha}^\dagger N \boldsymbol{\alpha}.
\]
Both $\boldsymbol{\alpha}^\dagger M \boldsymbol{\alpha}$ and $\boldsymbol{\alpha}^\dagger N \boldsymbol{\alpha}$ are real since $M$ and $N$ are Hermitian. Moreover, $\boldsymbol{\alpha}^\dagger N \boldsymbol{\alpha} = (\varphi, \varphi) = \|\varphi\|^2 > 0$ (assuming $\varphi \neq 0$). Therefore:
\[
\lambda = \frac{\boldsymbol{\alpha}^\dagger M \boldsymbol{\alpha}}{\boldsymbol{\alpha}^\dagger N \boldsymbol{\alpha}} \in \mathbb{R}.
\]

\subsection*{(d) Orthogonality of eigenvectors}

The eigenvectors are \emph{not} in general orthogonal in the usual sense, but they satisfy a \emph{generalized orthogonality} with respect to $N$. If $\boldsymbol{\alpha}^{(k)}$ and $\boldsymbol{\alpha}^{(l)}$ are eigenvectors with distinct eigenvalues $\lambda_k \neq \lambda_l$, then from $M\boldsymbol{\alpha}^{(k)} = \lambda_k N\boldsymbol{\alpha}^{(k)}$:
\[
{\boldsymbol{\alpha}^{(l)}}^\dagger M \boldsymbol{\alpha}^{(k)} = \lambda_k\,{\boldsymbol{\alpha}^{(l)}}^\dagger N \boldsymbol{\alpha}^{(k)}, \qquad {\boldsymbol{\alpha}^{(k)}}^\dagger M \boldsymbol{\alpha}^{(l)} = \lambda_l\,{\boldsymbol{\alpha}^{(k)}}^\dagger N \boldsymbol{\alpha}^{(l)}.
\]
Taking the conjugate of the second equation and using Hermiticity of $M$ and $N$:
\[
(\lambda_k - \lambda_l)\,{\boldsymbol{\alpha}^{(l)}}^\dagger N \boldsymbol{\alpha}^{(k)} = 0.
\]
Since $\lambda_k \neq \lambda_l$, we get ${\boldsymbol{\alpha}^{(l)}}^\dagger N \boldsymbol{\alpha}^{(k)} = 0$, which means the trial functions $\varphi^{(k)}$ and $\varphi^{(l)}$ are orthogonal in $L_2$:
\[
(\varphi^{(l)}, \varphi^{(k)}) = 0.
\]

\subsection*{(e) Relation to the largest eigenvalue of $A$}

The largest value of $\lambda$ obtained from $M\boldsymbol{\alpha} = \lambda N\boldsymbol{\alpha}$ provides a \emph{lower bound} for the largest eigenvalue of $A$.

By the variational principle (Theorem~9.2 and the Rayleigh--Ritz method), the largest eigenvalue of $A$ is:
\[
\lambda_{\max}(A) = \sup_{\|f\|=1} (f, Af).
\]
The trial function $\varphi = \sum a_n \phi_n$ explores only a $\nu$-dimensional subspace of $H$, so:
\[
\max_{\boldsymbol{\alpha}} \frac{\boldsymbol{\alpha}^\dagger M \boldsymbol{\alpha}}{\boldsymbol{\alpha}^\dagger N \boldsymbol{\alpha}} \le \lambda_{\max}(A).
\]
The largest $\lambda$ from the generalized eigenvalue problem is the best approximation to $\lambda_{\max}(A)$ within the chosen basis, and by choosing the $\{\phi_n\}$ intelligently, one can obtain an excellent approximation. \hfill $\blacksquare$

\end{document}
